% ============================================================================
%  Federated and meta-learned blood-glucose forecasting across sites
%  PLOS Digital Health submission (PLOS template v3.8 Apr 2026)
%
%  STATUS: Introduction (with prior work folded in, PLOS house style) complete.
%  Abstract, Author summary, Materials and methods, Results, Discussion are
%  marked TODO for later sessions.
%
%  NOTE: this uses the PLOS class options and the `cite` package (numeric
%  Vancouver citations). The bibliography style is `plos2015` (the genuine PLOS
%  Vancouver BibTeX style; the current template ships an identical file renamed
%  plos2025.bst). plos2015.bst is committed alongside this .tex so it compiles
%  as-is on Overleaf. The bib database is references.bib.
% ============================================================================
\documentclass[10pt,letterpaper]{article}
\usepackage[top=0.85in,left=2.75in,footskip=0.75in]{geometry}

\usepackage{amsmath,amssymb}
\usepackage{changepage}
\usepackage{textcomp,marvosym}
\usepackage{cite}
\usepackage{nameref,hyperref}
\usepackage[right]{lineno}
\usepackage[nopatch=eqnum]{microtype}
\DisableLigatures[f]{encoding = *, family = * }
\usepackage[table]{xcolor}
\usepackage{array}

\newcolumntype{+}{!{\vrule width 2pt}}
\newlength\savedwidth
\newcommand\thickcline[1]{%
  \noalign{\global\savedwidth\arrayrulewidth\global\arrayrulewidth 2pt}%
  \cline{#1}%
  \noalign{\vskip\arrayrulewidth}%
  \noalign{\global\arrayrulewidth\savedwidth}%
}
\newcommand\thickhline{\noalign{\global\savedwidth\arrayrulewidth\global\arrayrulewidth 2pt}%
\hline
\noalign{\global\arrayrulewidth\savedwidth}}

% Text layout
\raggedright
\setlength{\parindent}{0.5cm}
\textwidth 5.25in
\textheight 8.75in

\usepackage[aboveskip=1pt,labelfont=bf,labelsep=period,justification=raggedright,singlelinecheck=off]{caption}
\renewcommand{\figurename}{Fig}

\bibliographystyle{plos2015}

\makeatletter
\renewcommand{\@biblabel}[1]{\quad#1.}
\makeatother

% Header and Footer with logo
\usepackage{lastpage,fancyhdr,graphicx}
\usepackage{epstopdf}
\pagestyle{fancy}
\fancyhf{}
\rfoot{\thepage/\pageref{LastPage}}
\renewcommand{\headrulewidth}{0pt}
\renewcommand{\footrule}{\hrule height 2pt \vspace{2mm}}
\fancyheadoffset[L]{2.25in}
\fancyfootoffset[L]{2.25in}
\lfoot{\today}

\begin{document}
\vspace*{0.2in}

% Title must be 250 characters or less. Sentence case.
\begin{flushleft}
{\Large
\textbf{Federated and meta-learned blood-glucose forecasting across sites: a domain-generalisation benchmark and a cross-country deployment}
}
\newline
\\
Name1 Surname\textsuperscript{1*}
\\
\bigskip
\textbf{1} Affiliation Dept/Program/Center, Institution Name, City, State, Country
\\
\bigskip

* correspondingauthor@institute.edu

\end{flushleft}

% Please keep the abstract below 300 words
\section*{Abstract}
% TODO (later session): structured abstract <300 words summarising the cross-site
% FL system + the domain-generalisation benchmark and headline findings.
[To be written.]

% Author Summary: 150-200 words, first person.
\section*{Author summary}
% TODO (later session): plain-language summary for the PLOS Digital Health readership.
[To be written.]

\clearpage
\newgeometry{top=0.85in,left=1in,right=1in,footskip=0.75in}
\linenumbers

% ===========================================================================
\section*{Introduction}

Type~1 diabetes (T1D) affects an estimated 8.4~million people worldwide and its
incidence is rising in many regions~\cite{idf2021,gregory2022t1dindex}. People living with T1D must
continually balance insulin, food, and activity to keep blood glucose within a
narrow target range; sustained excursions above or below that range drive both
the acute dangers of hypoglycaemia and the long-term microvascular and
cardiovascular complications of hyperglycaemia. Continuous glucose monitoring
(CGM), which reports interstitial glucose roughly every five minutes, is now
central to diabetes care; it underpins clinical outcomes such as time in range,
for which international consensus targets have been defined~\cite{battelino2019tir}. The same dense data stream makes
short-horizon glucose \emph{forecasting} clinically valuable: a reliable
prediction of where glucose will be in 30--60 minutes can trigger preemptive
alerts, inform bolus and snack decisions, and close the loop in automated
insulin-delivery systems~\cite{vettoretti2020advanced,bekiari2018ap}.

Short-horizon CGM forecasting has progressed from autoregressive and classical
machine-learning models to deep sequence models. Recurrent architectures with
predictive variance estimation established strong baselines on benchmark
cohorts~\cite{martinsson2020}, and standardised datasets and challenges shaped
much of the early evaluation~\cite{marling2020ohio}. The
transformer~\cite{vaswani2017attention} and its long-horizon time-series
variants~\cite{zhou2021informer,wu2021autoformer} have since been adapted to
glucose, with personalised, uncertainty-aware transformer forecasters reporting
state-of-the-art accuracy~\cite{sergazinov2023gluformer,sergazinov2024glucobench}.
These models are accurate when trained and evaluated on the same cohort, but how
well they transfer across cohorts is an open question~\cite{ghimire2024generalize}:
glucose dynamics, demographics, and devices differ from site to site, and a model
can degrade on populations it did not see in training~\cite{farahmand2025glumind}.
A model that has only seen one cohort therefore has little reason to generalise to
the next; robust generalisation instead calls for training over large, diverse
populations spanning many sites.

Large, diverse training data is exactly what is hardest to assemble. CGM traces
are sensitive personal health records, fragmented across hospitals, clinical
trials, device manufacturers, and national registries, and governed by
data-protection regimes such as the GDPR and HIPAA that restrict moving raw
records between institutions, let alone across borders. The data that improve
generalisation and the data that exist are thus mismatched: large, heterogeneous
corpora are what is needed, but what is available is many small, governed silos.
Federated learning (FL) addresses this tension~\cite{mcmahan2017fedavg}: sites
train on their own data and exchange only model updates, so a shared model is
learned without any raw record leaving its institution. On this basis FL has been
widely adopted for collaborative medical machine
learning~\cite{rieke2020future,sheller2020federated,kaissis2020secure}. Its
central difficulty is statistical heterogeneity: plain federated averaging
degrades when client distributions diverge, motivating proximal
regularisation~\cite{li2020fedprox} and control-variate variance
reduction~\cite{karimireddy2020scaffold}. Most of this work, however, is evaluated
in simulation, where ``clients'' are partitions of one dataset on a single
machine.

Even with federation, real CGM cohorts differ in device generation, demographics,
treatment protocols, and glycaemic regime, and these differences widen across
countries. A federated model must therefore cope with substantial domain shift: it
should generalise to patients, and even whole cohorts, it never trained on (domain
generalisation), and/or specialise to each site (personalisation). Domain
generalisation seeks models that transfer to unseen
distributions~\cite{zhou2022dgsurvey}, and meta-learning for domain generalisation
(MLDG) simulates train/test domain shift within each update so the model learns to
generalise rather than memorise~\cite{li2018mldg}. Bringing these ideas into
federation is an active area surveyed by Li et al.~\cite{li2025fedDGsurvey},
spanning several families. Meta-learning approaches adapt the model-agnostic
meta-learning objective to federation, for communication-efficient
convergence~\cite{chen2018fedmeta} or for models that personalise to a new client
in a few gradient steps~\cite{fallah2020perfedavg}. Personalisation approaches keep
a per-client model: APFL learns an adaptive convex mixture of the global and a
local model~\cite{deng2020apfl}, while Ditto regularises each client's personal
model toward the global one with a tunable proximal strength~\cite{li2021ditto}.
Other approaches instead align client representations, via model-contrastive
objectives~\cite{li2021moon}, federated contrastive learning for multi-centre
medical imaging~\cite{fedcl2023}, or adversarial domain
alignment~\cite{flda2023,fedadv2022}. Our benchmark covers the first two
families---meta-learning (MLDG) and personalisation (APFL, Ditto)---alongside
standard federated averaging (FedAvg, FedProx) and single-cohort local training,
and leaves representation-alignment methods to future work.

\textbf{Yet these methods are rarely compared head-to-head for glucose
forecasting.} They are seldom benchmarked against one another on a common task, and
evaluations rarely separate the two regimes that matter clinically---accuracy on
patients from the training cohorts (held-in) versus accuracy on entirely unseen
cohorts (out of distribution, OOD)---even though a method can behave quite
differently in the two. Closing this gap---comparing federated strategies on one
glucose task under a shared protocol, with held-in and OOD performance reported
separately---is the main focus of this paper.

A second, more practical observation is that almost all federated glucose work is
simulated rather than deployed, with ``clients'' as partitions of one dataset on a
single machine. Running a federation for real additionally requires networked,
physically separate sites to exchange updates and to start from a common model;
production FL frameworks can do this but carry substantial dependencies and
assumptions that can be ill-suited to a small clinical consortium that simply wants
a handful of sites, one cohort each, to train together. We therefore also ask a
feasibility question: whether a lightweight cross-site federation for CGM
forecasting can be built and run across real machines with minimal tooling.

This paper makes one primary contribution and one secondary one. \textbf{Our
primary contribution is a controlled domain-generalisation benchmark of federated
strategies for transformer-based glucose forecasting.} Using a common
encoder--decoder transformer and four real T1D cohorts as federated clients, with
additional held-out cohorts as OOD test sites, we compare single-cohort local
training, FedAvg, FedProx, MLDG, and two personalised-FL methods (APFL and Ditto)
under a multi-seed protocol, reporting held-in and OOD accuracy separately at the
clinically actionable 30-minute horizon.

\textbf{As a secondary contribution, we show that such a federation can be run for
real, not only simulated.} We build, from scratch, a lightweight cross-site FL
system for CGM forecasting: a stateless aggregation server that holds the global
model but no data, and clients that each wrap one institution's cohort,
communicating over HTTP using only the Python standard library. The server
broadcasts the initial weights so every client starts from an identical model, and
the same protocol carries FedAvg, FedProx, and MLDG training. We validate it
end-to-end across two physically separate machines connected over the public
internet, and describe how the design extends to a geographically distributed,
multi-country deployment.

Taken together, the benchmark shows which federated strategy to run and how to
evaluate it honestly, while the system shows that such a federation can be run for
real across institutions and borders, not just simulated.

% ===========================================================================
\section*{Materials and methods}

\subsection*{Datasets and federation setup}

% REVIEW (citation check 2026-06-19): Table~\ref{tab:datasets} patient counts
% disagree with the source dataset papers -- HUPA-UCM is 22 here vs 25 in
% Hidalgo et al. 2024; T1D-UOM is 14 here vs 17 in Alsuhaymi et al. 2025 (21
% enrolled, 4 withdrew). These are presumably our usable post-filtering subsets;
% either reconcile the numbers or state the filtering criterion explicitly.
We use five publicly available type-1 diabetes (T1D) continuous glucose
monitoring (CGM) datasets. Four serve as the federated clients --- HUPA-UCM
\cite{hidalgo2024hupaucm}, ABC4D, ARISES, and T1D-UOM
\cite{alsuhaymi2025t1duom} --- with \emph{one dataset per client}. Each dataset
originates from a distinct study, population, and device, so a per-dataset
partition makes every client a genuinely different domain rather than an
arbitrary shard of a pooled corpus; this is what lets the federation function as
a test of cross-domain generalisation. The fifth dataset, ReplaceBG, is never
seen during training and is used only as a held-out, out-of-distribution (OOD)
test client to measure how well a federated model transfers to a cohort that
contributed no data to training. Within each client, we partition patients
into disjoint train, validation, and test groups, splitting strictly by
patient ID so that no patient contributes data to more than one group. Every
window in the validation and test sets therefore comes from a patient never
seen during training, and reported performance measures generalisation to
unseen patients rather than to unseen windows of patients already observed.

Two of the four client cohorts, HUPA-UCM and T1D-UOM, together with the held-out
ReplaceBG cohort, are obtained through MetaboNet
\cite{wolff2026metabonet}, a recently released resource that consolidates many
public T1D datasets into a single harmonised schema with a common 5-minute grid
and unified CGM, insulin, and carbohydrate fields. Drawing these cohorts from
MetaboNet ensures their variables are already aligned to a shared format; the
remaining two cohorts, ABC4D and ARISES, are converted into the same schema
from their original releases. We do not adopt MetaboNet's own train/test split,
which allows the same patient to appear in both partitions; instead we pool the
data and re-split strictly by patient.

Table~\ref{tab:datasets} summarises each cohort. Mean and standard deviation of
CGM are reported in mg/dL, and time-in-range (TIR) is the fraction of CGM
readings in the clinical 70--180 mg/dL band.

\begin{table}[!ht]
\centering
\captionsetup{justification=centering,singlelinecheck=on}
\caption{\bf Federated client cohorts.}
\label{tab:datasets}
\begin{tabular}{lrrrrr}
\thickhline
dataset & patients & total windows & mean CGM & sd CGM & TIR (\%) \\
        &          &               & (mg/dL)  & (mg/dL) & 70--180  \\
\hline
HUPA-UCM & 22 & 291,728 & 140.9 & 56.6 & 72.2 \\
ABC4D    & 25 & 562,043 & 155.1 & 63.4 & 65.5 \\
ARISES   & 12 & 126,248 & 160.5 & 63.7 & 63.7 \\
T1D-UOM  & 14 & 300,796 & 151.6 & 59.9 & 70.5 \\
\thickhline
\end{tabular}
\begin{center}
{\small ReplaceBG (OOD test client) is excluded from training and reported
separately. ARISES is both the highest-glucose and the smallest cohort.}
\end{center}
\end{table}

\subsection*{Preprocessing and windowing}

Every cohort is resampled to a uniform 5-minute grid. We form supervised
examples with a sliding window of 72 samples (6\,h) of CGM context to forecast
the next 12 samples (60\,min), advancing the window one sample at a time
(stride~1). Windows that would bridge a sensor gap longer than 6 samples
(30\,min) are discarded, so every example is contiguous. CGM is normalised with
a single \emph{global} mean and standard deviation ($154.04 \pm 61.00$ mg/dL)
shared across all clients rather than per-client statistics; a common
normalisation keeps the global model's input scale identical at every site,
which is what allows weights aggregated across clients to be applied directly to
the OOD cohort. Alongside CGM, each example carries a time-of-day mark encoded
as $(\sin 2\pi\,\mathrm{frac},\ \cos 2\pi\,\mathrm{frac})$, where $\mathrm{frac}$
is the fraction of the day, to expose circadian structure.

\subsection*{Forecasting model}

The forecaster is a sequence-to-sequence (encoder--decoder) Transformer that maps
the 72-step CGM context to the 12-step horizon. It uses
$d_{\text{model}}=128$, 8 attention heads, 6 encoder layers, 2 decoder layers,
feed-forward width $d_{\text{ff}}=2048$, ReLU activations, and dropout~$0.05$
($\approx 4.9$\,M parameters). The decoder predicts a single CGM channel
($c_{\text{out}}=1$) and the model is trained with mean-squared-error (MSE) loss.
The same architecture is used for every client and every federated strategy, so
all comparisons differ only in how the model is trained, not in capacity.

\paragraph{Bolus input variant}

Not quite sure if I should include this. It is somewhat irrelevant?

To test whether insulin information improves forecasting, we additionally train a
two-channel variant ($\texttt{enc\_in}=2$) in which a second encoder channel
carries $\log(1+\mathrm{bolus}_{\mathrm{IU}})$, the log-transformed bolus insulin
dose; bolus events are sparse (roughly 1\% of timesteps are non-zero). The
decoder remains CGM-only. We report this variant as a delta against the
CGM-only model.

\subsection*{Federated learning system}

TO DO: Add details about where training takes place physically.

Training is coordinated by a from-scratch, networked federated learning system
in which one institution runs a server and each cohort runs a client. The server
holds the global model weights but never sees any data: at round one it mints a
single seeded set of initial weights and broadcasts them, so every client begins
from byte-identical parameters regardless of hardware. Each round, clients fetch
the current global weights, train locally on their own cohort, and return updated
weights; the server aggregates them by FedAvg and advances to the next round.
Clients communicate with the server over plain HTTP through a small task
state-machine (train\,$\rightarrow$\,validate\,$\rightarrow$\,test), so the same
protocol carries FedAvg, FedProx, and MLDG without modification --- the choice of
strategy is entirely client-side, and the server always performs the same FedAvg
aggregation. We validate the system end-to-end across two physically separate
machines communicating over the public internet, and describe how the design
generalises to a geographically distributed, multi-country deployment.

\subsection*{Federated strategies and baselines}

\textbf{Local (baseline).} Our primary comparator is local training: a separate
model trained on a single cohort's data only, with no federation. The central
question of the benchmark is whether federation improves on this baseline ---
both when evaluated on the same cohorts that participated (held-in) and on the
unseen OOD cohort. For OOD we report the mean over the four single-cohort models.

\textbf{FedAvg / FedProx.} Standard federated averaging, and FedProx, which adds
a proximal term $(\mu/2)\lVert w - w_{\text{global}}\rVert^2$ to each client's
objective ($\mu=0.05$) to limit client drift.

\textbf{MLDG.} Meta-learning for domain generalisation, adapted to the federated
setting: at each client step the patient domains are split into meta-train and
meta-test partitions, and the model is updated with a second-order meta-gradient
(the inner step reuses the outer Adam optimiser state). MLDG and FedProx alter
only client-side training; the server still FedAvg-aggregates.

\textbf{Personalised FL (PFL).} APFL learns a per-client mixture
$\alpha v_i + (1-\alpha)w$ of a personal model $v_i$ and the global model $w$
with an adaptively updated $\alpha$ (initial $0.25$, $\alpha$ learning
rate~$0.1$); we also include a decoupled APFL variant. Ditto trains a personal
model pulled toward the global one by a proximal term ($\mu\in\{0.01,0.1,1.0\}$).
These methods produce one personalised model per client and \emph{no} single
global model, so --- like the local baseline's per-cohort models --- they are
evaluated held-in only and cannot be applied to the OOD cohort.

\subsection*{Evaluation protocol and metrics}

We evaluate every strategy in two regimes. \emph{Held-in} performance is measured
on each participating cohort's own held-out patients, isolating predictive
accuracy at the contributing sites. \emph{OOD} performance is measured on
ReplaceBG, which contributed no training data, isolating cross-domain transfer;
only strategies that yield a single deployable global model (Local-mean, FedAvg,
FedProx, MLDG) receive an OOD score. We report root-mean-squared error (RMSE) in
mg/dL at the 30-minute and 60-minute horizons. All runs use early stopping on
average validation MSE (patience~5 over a maximum of 25 communication rounds),
and every configuration is repeated over 5 random seeds. We
report the mean and standard deviation across seeds.

\subsection*{Implementation and training details}

All models are trained with the Adam optimiser (learning rate~$10^{-4}$, weight
decay~$10^{-3}$, batch size~256, one local epoch per communication round). The
federated system is implemented in Python using only the standard library for
networking.

% TODO: data availability + ethics statement (public de-identified datasets);
% architecture figure for the FL system.

\section*{Results}

Table~\ref{tab:main} reports held-in RMSE on each client's unseen test patients
and OOD RMSE on ReplaceBG, at the 30-minute horizon, for the CGM-only model.
% TODO: add 60-minute-horizon numbers once the corresponding metrics are exported.

\begin{table}[!ht]
\centering
\captionsetup{justification=centering,singlelinecheck=on}
\caption{\bf CGM-only RMSE (mg/dL) at the 30-minute horizon.}
\label{tab:main}
\begin{tabular}{lrrrrrr}
\thickhline
 & \multicolumn{5}{c}{held-in (per client)} & OOD \\
\cline{2-6}
strategy & HUPA-UCM & ABC4D & ARISES & T1D-UOM & avg & ReplaceBG \\
\hline
Local (per cohort) & 18.81 & 19.95 & 22.04 & 19.80 & 20.15 & 23.54 \\
\hline
FedAvg  & 18.64 & 19.80 & 22.38 & 19.53 & 20.09 & 21.43 \\
FedProx & 18.89 & 19.94 & 22.60 & 19.67 & 20.28 & 21.57 \\
MLDG    & \textbf{18.46} & \textbf{19.68} & 22.30 & 19.53 & \textbf{19.99} & \textbf{21.34} \\
\hline
APFL          & 19.11 & 19.88 & 22.53 & 19.61 & 20.28 & --- \\
APFL-decoupled & 19.37 & 19.90 & 22.47 & 19.73 & 20.37 & --- \\
Ditto ($\mu{=}0.01$) & 18.90 & 19.90 & \textbf{21.98} & 19.64 & 20.11 & --- \\
Ditto ($\mu{=}0.1$)  & 18.90 & 19.93 & 22.28 & 19.66 & 20.19 & --- \\
Ditto ($\mu{=}1.0$)  & 18.74 & 19.77 & 22.35 & \textbf{19.50} & 20.09 & --- \\
\thickhline
\end{tabular}
\begin{center}
{\small Mean over 5 random seeds (4 where a run is missing). \emph{avg} is the
mean of the four held-in client means. Local held-in is each cohort's own
single-cohort model; Local OOD (23.54) is the mean of the four single-cohort
models on ReplaceBG. Personalised strategies (APFL, Ditto) yield no single
global model and so receive no OOD score. Best value per column in bold.}
\end{center}
\end{table}

\paragraph{Federation matches local accuracy held-in and transfers far better
out-of-distribution.} On the contributing cohorts, the global FedAvg and MLDG
models are on par with cohort-specific local training (held-in average $20.09$
and $19.99$ vs.\ $20.15$ mg/dL): pooling four heterogeneous cohorts through
federation does not cost in-domain accuracy despite each client seeing only the
shared global weights. The decisive difference is out-of-distribution. On the
never-seen ReplaceBG cohort, FedAvg ($21.43$) and MLDG ($21.34$) improve on the
mean single-cohort model ($23.54$) by more than $2$ mg/dL, and they do so far
more reliably: the local baseline's OOD error is both higher and much noisier
across seeds (sd $1.65$ mg/dL), driven by the ARISES-only model, which collapses
on ReplaceBG ($28.51 \pm 6.85$). A model trained on a single cohort can transfer
catastrophically; the federated model never does.

\paragraph{MLDG is the strongest global strategy.} Among the strategies that
produce a deployable global model, MLDG is best on the held-in average and best
OOD, and is the top method on two of the four clients (HUPA-UCM, ABC4D). FedProx
is consistently the weakest of the three global strategies here, both held-in and
OOD. The personalised methods do not change the held-in picture --- their
averages cluster around FedAvg's ($20.1$--$20.4$) --- and by construction they
cannot be deployed on the OOD cohort.

\paragraph{Finetuning with replaceBG} 
So FL + finetuning works better on replace BG. Experiment still running so will report on that later.

\section*{Discussion}

\paragraph{Principal findings.} Across four real T1D cohorts federated as
distinct domains, with a fifth never-seen cohort as an out-of-distribution (OOD)
test, federation matched single-cohort local training in-domain and transferred
substantially better out of it. Held-in, the global FedAvg and MLDG models were
on par with cohort-specific local training (average RMSE $20.09$ and $19.99$
vs.\ $20.15$ mg/dL at the 30-minute horizon); OOD, both improved on the mean
single-cohort model (ReplaceBG $21.43$ and $21.34$ vs.\ $23.54$ mg/dL) and did
so far more reliably. Among the strategies that yield a single deployable global
model, MLDG was the strongest both held-in and OOD. Beyond the benchmark, we
showed that the federation runs for real: the same protocol trains FedAvg,
FedProx, and MLDG across physically separate machines communicating over the
public internet, with no raw data leaving any site.

\paragraph{Federation as a robustness guarantee, not only an accuracy gain.} The
clearest message of the benchmark is not that federation is marginally more
accurate in-domain --- the held-in gap over local training is small --- but that
it removes a deployment risk that single-cohort training cannot. A model trained
on one cohort can transfer catastrophically: the ARISES-only model, accurate on
its own patients ($22.04$ mg/dL), collapsed on ReplaceBG ($28.51 \pm 6.85$ mg/dL)
and drove most of the local baseline's high OOD variance (sd $1.65$ mg/dL). The
difficulty is that one cannot know in advance which site's model will be the one
that fails. Federation removes this gamble: every global model we trained
generalised to the unseen cohort, and none collapsed. For a clinical setting in
which distribution shift is inevitable --- new patients, new device generations,
changing treatment protocols --- this robustness is itself the deployment
argument, independent of any in-domain accuracy difference.

\paragraph{Generalisation versus personalisation, and how to have both.} Our
results expose a genuine tension. The personalised strategies (Ditto, APFL) and
dedicated single-cohort models were the \emph{only} approaches that beat the
federated models on the hardest and smallest in-domain cohort: on ARISES, Ditto
($\mu{=}0.01$, $21.98$ mg/dL) and the dedicated ARISES model ($22.04$ mg/dL) edged
MLDG ($22.30$). Yet these same methods produce no single global model and so
cannot serve the OOD cohort at all --- exactly the regime in which the global
models are most valuable. The two objectives pull in opposite directions:
specialising to a site helps that site but forfeits transfer; generalising across
sites transfers but leaves per-site accuracy on the table. The way to obtain both
is to treat the generalised global model as a foundation and adapt it locally,
rather than choosing one regime outright.

\paragraph{Finetuning a federated model for a new site.}
% TODO (finetuning experiment in progress): report finetuning the global federated
% model on ReplaceBG vs.\ training a model on ReplaceBG from scratch. Expected
% message: the federated model is not only a better zero-shot OOD predictor but
% also a better \emph{starting point} --- finetuning it on a new site's data beats
% local-only training there, so a newly joining centre benefits twice (a robust
% global model out of the box, then a personalised model that exceeds what its own
% data could train alone). Fill in numbers and the delta over local once the runs
% complete.
[Finetuning results to be added once the experiment completes.]

\paragraph{Why meta-learning helps.} That MLDG was the best global strategy, both
held-in on average and OOD, is consistent with its objective: by splitting each
client's patient domains into meta-train and meta-test partitions and updating so
that progress on one transfers to the other, MLDG optimises directly for the
cross-domain transfer we evaluate, rather than for in-distribution fit alone. The
effect is real but modest --- MLDG's margin over plain FedAvg is small relative to
the gap between any federated method and single-cohort training. The dominant
effect is federation itself; the choice of federated strategy is a second-order
refinement, with MLDG the best-supported default and FedProx the weakest of the
three here.

\paragraph{Implications for a clinical consortium.} These findings give a small
group of institutions a concrete reason to federate and a recipe for doing so.
A participating centre loses nothing in-domain, gains a model that will not
collapse under distribution shift, and --- once finetuning is accounted for ---
can adapt the shared model to outperform what its own data could train in
isolation. Because no raw record leaves any institution, this collaboration is
compatible with data-protection regimes that would otherwise forbid pooling, and
our cross-machine deployment shows the necessary tooling is light enough for a
consortium without dedicated federated-learning infrastructure. The benefit also
compounds: each additional cohort makes the shared model more representative, so
the incentive to join strengthens as the federation grows.

\paragraph{Strengths.} Three features distinguish this study from prior federated
glucose work. First, it compares local training, FedAvg, FedProx, MLDG, and two
personalised-FL families head-to-head on one task under a shared protocol, rather
than reporting a single method in isolation. Second, it separates held-in from
OOD performance, a distinction that prior evaluations rarely make explicit but
that proves decisive: the strategies are nearly indistinguishable in-domain yet
differ sharply in transfer. Third, the cohorts are partitioned strictly by
patient and normalised with a single shared statistic, and the federation is run
across real, physically separate machines rather than simulated as shards of one
dataset on one host.

\paragraph{Limitations.} Several constraints bound these conclusions. Our
strongest claim --- robust OOD transfer --- rests on a single held-out cohort
(ReplaceBG); a larger panel of OOD cohorts is needed to establish how broadly the
generalisation holds. A small number of seeds are missing (the OOD run for FedAvg
and FedProx at one seed; one T1D-UOM seed), and we report seed means and standard
deviations without a formal significance test, so the smaller margins --- notably
MLDG over FedAvg --- should be read as suggestive rather than definitive. We
report the 30-minute horizon here and leave the 60-minute horizon to a fuller
analysis. Privacy in our system means that raw data never moves, but we provide no
formal guarantee: we do not add secure aggregation or differential privacy, and a
determined adversary could in principle attack the exchanged updates. Finally, all
five cohorts are adult T1D populations, leaving paediatric and type-2 transfer
untested, and we deliberately scope the benchmark to meta-learning and
personalisation, leaving representation-alignment federated methods to future
work.

\paragraph{Future work.} The natural next steps follow directly from these
limitations: evaluating against more OOD cohorts and, ultimately, a true
multi-country deployment; adding secure aggregation and differential privacy to
turn the no-data-movement property into a formal guarantee; extending to longer
horizons and to event-aware inputs such as insulin and carbohydrate information;
and systematically combining the generalisation of MLDG with the per-site
adaptation of finetuning, so that the global model and its local specialisations
are trained jointly rather than in sequence. As populations and devices drift over
time, continual and adaptive federation --- updating the shared model as cohorts
evolve --- is a further direction this framework is positioned to support.

% ###########################################################################
% #  TEMPORARY SECTION -- NOT FOR SUBMISSION                                 #
% #  Scratch summary of prior published results on our four cohorts, to see  #
% #  where they sit relative to our numbers (Table~\ref{tab:main}).          #
% #  Plain-text citations (these refs are not yet in references.bib).        #
% #  DELETE before submission, or migrate the comparable rows into Results/  #
% #  Discussion with proper \cite keys.                                      #
% ###########################################################################
\clearpage
\section*{[TEMP] Prior published results on our cohorts}

\textbf{Read the protocol column before comparing.} Our benchmark uses a strict
\emph{by-patient} split, a 5-minute grid, and reports RMSE in mg/dL at the
30-minute horizon as mean\,$\pm$\,sd \emph{across 5 seeds}. Almost every paper
below differs on at least one of these axes --- most use a \emph{by-time}
(chronological within-patient) split, which is an easier task because the test
patients are seen during training, and several report sd \emph{across patients}
rather than seeds. Numbers are therefore \emph{indicative}, not head-to-head.

\begin{table}[!ht]
\centering
\captionsetup{justification=centering,singlelinecheck=on}
\caption{\bf [TEMP] External results on HUPA-UCM, ABC4D, ARISES, T1D-UOM.}
\label{tab:prior}
\begin{tabular}{lllll}
\thickhline
study (best model) & cohort & RMSE@30 (mg/dL) & split & sd over \\
\hline
\multicolumn{5}{l}{\emph{Comparable horizon/metric (RMSE @ 30\,min)}} \\
GlucoFM-Bench (LSTM, full-shot) & HUPA-UCM & 16.71 (4.82) & by-time 80/20 & patients \\
GlucoFM-Bench (LSTM, full-shot) & T1D-UOM  & 20.81 (4.26) & by-time 80/20 & patients \\
FCNN (Zhu, evidential+meta)     & ARISES   & 20.23 $\pm$ 3.38 & by-time 64/16/20 & patients \\
FCNN (Zhu, evidential+meta)     & ABC4D    & 20.25 $\pm$ 2.60 & by-time 64/16/20 & patients \\
\hline
\multicolumn{5}{l}{\emph{Not directly comparable (see notes)}} \\
Tan \& McBeth (Transformer-evid.) & HUPA-UCM & MARD 4.14\% & by-time 60/20/20 & --- \\
AEGIS (XGBoost)                   & T1D-UOM  & 0.52 mmol/L @15\,min & by-patient (LOPO) & --- \\
Manchanda (aggregated LSTM)       & HUPA-UCM & 20.50 (5.66)$^\dagger$ & by-time 60/20/20 & patients \\
Kolev (LightGBM-SDE)              & HUPA-UCM$^\ddagger$ & 28.74 $\pm$ 9.4 & by-time 80/20 & patients \\
\thickhline
\end{tabular}
\begin{center}
{\small Our by-patient RMSE@30 for reference (Table~\ref{tab:main}, best global
strategy MLDG): HUPA-UCM 18.46, ABC4D 19.68, ARISES 22.30, T1D-UOM 19.53 mg/dL.
$^\dagger$Manchanda's figure is a 0--60\,min \emph{window-averaged} RMSE (12-step
output), not a 30- or 60-min point error, so it is not a horizon-specific number.
$^\ddagger$Kolev ran on \emph{synthetic} Bergman-model data, not the real
HUPA-UCM recordings, and pooled it with a synthetic D1NAMO cohort.}
\end{center}
\end{table}

\paragraph{GlucoFM-Bench (Lu et al., 2026) --- HUPA-UCM, T1D-UOM.} A benchmark of
time-series foundation models (Chronos-2, Moirai~2.0, Timer, TimesFM~2.5, plus
LSTM/Transformer baselines and TimeLLM/CALF) for \emph{univariate} CGM
forecasting across 15 datasets. \emph{Split:} chronological per-participant
(first 80\% train / last 20\% test), explicitly chosen to avoid the variance of
participant-level splits --- i.e.\ the opposite choice to ours. \emph{Features:}
CGM only, no insulin/carbs/covariates, 12\,h context, 15-min grid.
\emph{Best RMSE@30 (full-shot):} a lightweight LSTM, HUPA-UCM 16.71 (4.82),
T1D-UOM 20.81 (4.26) mg/dL; sd is across patients. These are the most directly
relevant external numbers, but the by-time split makes them an easier task than
ours; their HUPA-UCM 16.7 vs our by-patient 18.5 is consistent with that.

\paragraph{FCNN (Zhu et al., IEEE TBME 2023) --- ARISES, ABC4D.} A personalised
forecaster: bidirectional-GRU + attention, with evidential regression for
uncertainty and MAML/Reptile meta-learning for fast per-subject adaptation;
baselines SVR, RFR, ARIMA, Bi-LSTM, Transformer, CRNN. \emph{Split:} two-step
chronological per-subject (64/16/20). \emph{Features:} CGM + meal carbohydrate +
bolus insulin (5-min). \emph{Best RMSE@30 (FCNN):} ARISES 20.23\,$\pm$\,3.38,
ABC4D 20.25\,$\pm$\,2.60 mg/dL; sd across subjects, single run (no seed
variation). Note this is the same ARISES/ABC4D source (Imperial College) we use;
our by-patient ARISES (MLDG 22.30, Ditto 21.98) and FCNN's by-time 20.23 differ
in the expected direction given the harder split and the extra insulin/carb
inputs they use.

\paragraph{Tan \& McBeth (2026) --- HUPA-UCM, MARD only.} Uncertainty-aware
LSTM/GRU/Transformer with MC-dropout and deep evidential regression. \emph{Split:}
chronological per-patient 60/20/20. \emph{Features:} CGM + bolus + carbs + one of
\{heart rate, steps, calories, basal\}. Reports \emph{no RMSE} --- accuracy is
MARD and DTS error-grid zone-A; best model (Transformer-evidential) reaches MARD
4.14\% and zone-A 96.8\% at 30\,min with heart-rate input. Useful only as a
qualitative MARD point.

\paragraph{AEGIS (Pansheriya \& Asadinia, IEEE CCWC 2026) --- T1D-UOM.} An
end-to-end prediction + risk-classification + bolus-advice pipeline; prediction by
MLP vs XGBoost (XGBoost best). \emph{Split:} \emph{leave-one-patient-out} --- the
only external paper that, like us, evaluates on unseen patients. But it forecasts
only 10/15\,min (explicitly abandoned $>$30\,min for ``inadequate data density''),
reports in mmol/L, and gives no sd. XGBoost @15\,min RMSE 0.517 mmol/L
($\approx$9.3 mg/dL). \emph{Features:} CGM + bolus + carbs + activity (MET/steps)
+ circadian, 14 engineered features. Same by-patient philosophy as ours but a much
shorter horizon, so not directly comparable on RMSE@30.

\paragraph{Lower-value comparisons.} \emph{Manchanda et al.\ (Diabetology 2025)}
trains individualised vs aggregated LSTMs on HUPA-UCM (CGM + carbs + bolus + basal,
chronological 60/20/20); its headline 20.50 (5.66) mg/dL is a 0--60\,min
window-averaged RMSE over a 12-step output, \emph{not} a horizon-specific number,
so it should not be cited as a 30- or 60-min result. \emph{Kolev et al.\ (Appl.\
Sci.\ 2026)} benchmarks 10 ML methods (best: a hybrid LightGBM + stochastic Bergman
model, 28.74\,$\pm$\,9.4 mg/dL @30\,min) but ran entirely on \emph{synthetic} data
(the real HUPA-UCM server returned HTTP~403) pooled with synthetic D1NAMO, so it
does not constitute a real-data baseline.

\paragraph{Takeaway for positioning.} On HUPA-UCM and T1D-UOM the only strong
real-data RMSE@30 baselines are GlucoFM-Bench's, and on ARISES/ABC4D they are
FCNN's --- but all four use by-time splits, so our by-patient numbers are not
expected to beat them and should not be framed as doing so. The honest framing is
that our results are in the same 17--22 mg/dL range as prior work \emph{under a
strictly harder (by-patient) protocol}, and that no prior paper reports all four
cohorts, separates held-in from OOD, or evaluates across-patient generalisation as
we do.

% ###########################################################################
% #  END TEMPORARY SECTION                                                   #
% ###########################################################################

% ===========================================================================
% Compile references.bib with the plos2025.bst style (ships with PLOS template).
\bibliography{references}

\end{document}
